\section{Sums and products in \texorpdfstring{$\N$}{N} and \texorpdfstring{$\Q$}{Q}}
\label{SP:sec}

This section is about three questions from~\cite[\S5.1]{kra2023problems} that ask for sums and
products of infinite sets at the same time.
In \cref{SP:sec:u51a} we answer Question~\S5.1(a) positively: the two explicit sets named
in~\cite[\S5.1]{kra2023problems}, namely Beatty sequences and the set of integers with an even
number of prime factors, both contain $B+C$ and $BC$ for some infinite $B,C\subset\N$.
In \cref{SP:sec:q54} we answer Question~5.4 of~\cite{kra2023problems} negatively. The counterexample
is a set built at logarithmic scale; it has logarithmic density $3/10$ and it is aperiodic.
The same set answers Question~5.2 of~\cite{kra2023problems} negatively if ``positive density''
is read as positive upper density (\cref{SP:rem:q52}).
In \cref{SP:sec:q56} we answer Question~5.6 of~\cite{kra2023problems} negatively, as it is literally
stated. The counterexample comes from the sign obstruction at the real place, which survives the
removal of a set of density zero.

\smallskip
\noindent\emph{Notation.}
For $x\in\R$ we write $\{x\}\in[0,1)$ for the fractional part of $x$ and $\|x\|$ for the distance
from $x$ to the nearest integer.
We identify $\T=\R/\Z$ with $[0,1)$, and we write $\|x-y\|$ for the distance in $\T$.
For $X,Y\subset\N$ (or $\Q$) we write $X+Y=\{x+y:x\in X,\,y\in Y\}$, $XY=\{xy:x\in X,\,y\in Y\}$,
$X+t=\{x+t:x\in X\}$ and $dX=\{dx:x\in X\}$.
Recall from~\cite[eq.~(19)]{kra2023problems} that $A\subset\N$ is \emph{aperiodic} if
\begin{equation}\label{SP:eq:aperiodic}
\lim_{N\to\infty}\frac1N\sum_{n=1}^N\one_A(n)\,\e(\alpha n)=0\qquad\text{for all }\alpha\in\Q\setminus\Z,
\end{equation}
where $\e(x)=e^{2\pi i x}$.

%-----------------------------------------------------------------------------------------------
\subsection{The explicit examples}
\label{SP:sec:u51a}

Question~5.2 of~\cite{kra2023problems} asks whether every aperiodic $A\subset\N$ of positive
density contains $B+C$ and $BC$ for some infinite sets $B,C\subset\N$. Just before that question
the survey names two typical aperiodic sets: the set of all numbers with an even number of prime
factors, and the Beatty sequences $\{\lfloor n\alpha+\gamma\rfloor:n\in\N\}$, where $\alpha>0$ is
irrational and $\gamma\in\R$. It then remarks: ``Even for the explicit examples we have mentioned,
we consider the answer to be interesting.''

\begin{question}{\S5.1(a)}\label{SP:q:u51a}
Let $A$ be either a Beatty sequence $\{\lfloor n\alpha+\gamma\rfloor:n\in\N\}$, with $\alpha>0$
irrational and $\gamma\in\R$, or the set of all $n\in\N$ with an even number of prime factors. Are there infinite sets
$B,C\subset\N$ with $B+C\subset A$ and $BC\subset A$?
\end{question}

The survey does not say whether prime factors are counted with multiplicity. We answer the
question for both readings. The answer is yes in all cases (\cref{SP:cor:beatty,SP:thm:liouville}).

\subsubsection{Bohr sets and Beatty sequences}

\begin{theorem}\label{SP:thm:bohr}
Let $\vartheta\in\R\setminus\Q$, $\delta\in\R$, let $J\subset\T$ be a nonempty open set and let
$M_0\in\N$. Suppose that
\[
A\supseteq\{m\in\N:\ m\ge M_0,\ m\vartheta+\delta\bmod 1\in J\}.
\]
Then there is an infinite set $Z\subset\N$ such that $z+z'\in A$ and $zz'\in A$ for all distinct
$z,z'\in Z$. In particular, there are disjoint infinite sets $B,C\subset\N$ with $B+C\subset A$
and $BC\subset A$.
\end{theorem}

\begin{proof}
Since $J$ is open and nonempty, there are $\xi_0\in\T$ and $\ell\in(0,1)$ such that every
$x\in\T$ with $\|x-\xi_0\|<\ell/2$ lies in $J$. Fix $q\in\N$ with $q\ge 4/\ell$.
The points $2r/q$ ($r\in\Z$) form a subgroup of $\T$ in which consecutive points are at distance
at most $2/q$. Hence we can choose $r\in\Z$ with $\|2r/q+\delta-\xi_0\|\le 1/q\le\ell/4$.
Put $\xi=2r/q+\delta$. Then
\begin{equation}\label{SP:eq:xiJ}
\|x-\xi\|<\ell/4\ \Longrightarrow\ \|x-\xi_0\|<\ell/2\ \Longrightarrow\ x\in J .
\end{equation}

We choose $z_1<z_2<\cdots$ recursively. Put $z_0=1$, and suppose $z_0,\dots,z_{n-1}$ have been
chosen. Let $\eps_n=\ell/(8z_{n-1})$. We look for $z_n$ of the form $z_n=qk+2$ with $k\in\N$.
Then $z_n\vartheta-r/q=k(q\vartheta)+(2\vartheta-r/q)$. Since $q\vartheta$ is irrational, the
sequence $(kq\vartheta\bmod1)_{k>K}$ is dense in $\T$ for every $K$. So we can choose $k$ such that
$z_n=qk+2>\max(z_{n-1},M_0)$ and $\|z_n\vartheta-r/q\|<\eps_n$.
Thus every $n\ge1$ satisfies
\[
z_n\equiv 2\pmod q,\qquad z_n\vartheta=K_n+\frac rq+e_n,\quad K_n\in\Z,\ |e_n|<\eps_n=\frac{\ell}{8z_{n-1}}\le\frac\ell8 .
\]
Let $1\le i<j$. Write $z_i=qk_i+2$ with $k_i\in\Z$. We compute the two phases modulo $1$.
\begin{itemize}
\item Sums: $(z_i+z_j)\vartheta\equiv 2r/q+e_i+e_j$, and $|e_i+e_j|<\ell/4$.
\item Products: $z_iz_j\vartheta=z_iK_j+z_ir/q+z_ie_j$. Here $z_iK_j\in\Z$ and
$z_ir/q=k_ir+2r/q\equiv 2r/q$. Moreover $|z_ie_j|<z_i\ell/(8z_{j-1})\le\ell/8$, because
$z_i\le z_{j-1}$. So $z_iz_j\vartheta\equiv 2r/q+z_ie_j$.
\end{itemize}
Hence both $(z_i+z_j)\vartheta+\delta$ and $z_iz_j\vartheta+\delta$ lie within distance
$\ell/4$ of $\xi$ in $\T$, so they lie in $J$ by \eqref{SP:eq:xiJ}. Since
$z_i+z_j$ and $z_iz_j$ are at least $M_0$, both lie in $A$.
Take $Z=\{z_n:n\ge1\}$, $B=\{z_1,z_3,z_5,\dots\}$ and $C=\{z_2,z_4,z_6,\dots\}$.
\end{proof}

\begin{corollary}\label{SP:cor:beatty}
Let $\alpha>0$ be irrational, let $\gamma\in\R$, and let
$A=\{\lfloor n\alpha+\gamma\rfloor:n\in\N\}$. Then there are infinite sets $B,C\subset\N$ with
$B+C\subset A$ and $BC\subset A$.
\end{corollary}

\begin{proof}
Suppose first that $\alpha<1$. Then consecutive terms of the nondecreasing sequence
$\lfloor n\alpha+\gamma\rfloor$ differ by $0$ or $1$, and the sequence tends to infinity. So $A$
contains every integer $m\ge\lfloor\alpha+\gamma\rfloor$, and any infinite sets $B,C$ of
sufficiently large integers work.

Now suppose $\alpha>1$ and put $\beta=1/\alpha\in(0,1)$. An integer $m$ lies in $A$ if and only
if some $n\in\N$ satisfies $m\le n\alpha+\gamma<m+1$, that is, $y_m\le n<y_m+\beta$, where
$y_m=(m-\gamma)\beta$. The least integer that is at least $y$ is $y+\{-y\}$. So the interval
$[y_m,y_m+\beta)$ contains an integer if and only if $\{-y_m\}<\beta$. When $m>\gamma$, this
integer is at least $y_m>0$, so it lies in $\N$. Therefore, for $m>\gamma$,
\[
m\in A\iff \{m\vartheta+\delta\}\in[0,\beta),\qquad\text{where }\vartheta=-\beta,\ \delta=\gamma\beta .
\]
Here $\vartheta$ is irrational. Apply \cref{SP:thm:bohr} with $J=(0,\beta)$ and any
$M_0>\gamma$.
\end{proof}

\begin{remark}\label{SP:rem:bohr}
\Cref{SP:thm:bohr} does not use any density assumption. It applies to every set that contains an
inhomogeneous Bohr set with an irrational frequency, up to finitely many elements; this includes
all Beatty sequences. The argument is elementary, and we expect the statement to be folklore. We
did not find it in the literature.
\end{remark}

\subsubsection{The Liouville function}

Let $\Omega(n)$ and $\omega(n)$ be the numbers of prime factors of $n$ counted with and
without multiplicity, respectively. Let $\lambda(n)=(-1)^{\Omega(n)}$ be the Liouville function.
It is completely multiplicative: $\lambda(mn)=\lambda(m)\lambda(n)$ for all $m,n\in\N$. Put
\[
L_+=\{n\in\N:\Omega(n)\text{ is even}\}=\{n:\lambda(n)=1\},\qquad
L'_+=\{n\in\N:\omega(n)\text{ is even}\}.
\]

\begin{theorem}\label{SP:thm:liouville}
For each $E\in\{L_+,L'_+\}$ there are infinite sets $B,C\subset\N$ with $B+C\subset E$ and
$BC\subset E$.
\end{theorem}

\begin{proof}
Let $S$ be the set of primes $p\equiv5\pmod6$. There are infinitely many such primes: if
$p_1,\dots,p_r$ were all of them, then $6p_1\cdots p_r-1\equiv5\pmod6$ would be coprime
to $6p_1\cdots p_r$; so all its prime factors would be $\equiv1\pmod6$, which is impossible. Let $E\in\{L_+,L'_+\}$, and colour each
two-element subset $\{s,s'\}$ of $S$ according to whether $s+s'\in E$. By the infinite Ramsey
theorem for pairs~\cite{ramsey1930problem} there is an infinite set $S'\subset S$ such that either
$s+s'\in E$ for all distinct $s,s'\in S'$, or $s+s'\notin E$ for all distinct $s,s'\in S'$. Split
$S'$ into two disjoint infinite sets $B_0$ and $C_0$, and let $b_0\in B_0$, $c_0\in C_0$. Then
$b_0\neq c_0$ are primes different from $2$ and $3$, and $b_0+c_0$ is even and $\equiv1\pmod3$.
Also $\Omega(b_0c_0)=\omega(b_0c_0)=2$, so $b_0c_0\in L_+\cap L'_+$.

If $s+s'\in E$ for all distinct $s,s'\in S'$, then $B=B_0$ and $C=C_0$ work. Otherwise we put
$B=dB_0$ and $C=dC_0$, where $d=2$ if $E=L_+$ and $d=6$ if $E=L'_+$.
\begin{itemize}
\item $E=L_+$: $\Omega(2(b_0+c_0))=\Omega(b_0+c_0)+1$ is even, and
$\Omega(4b_0c_0)=4$.
\item $E=L'_+$: since $2\mid b_0+c_0$ and $3\nmid b_0+c_0$, we get
$\omega(6(b_0+c_0))=\omega(b_0+c_0)+1$, which is even. Since $b_0,c_0\notin\{2,3\}$,
$\omega(36b_0c_0)=4$.
\end{itemize}
In both cases $b+c\in E$ and $bc\in E$ for all $b=db_0\in B$ and $c=dc_0\in C$.
\end{proof}

For $L_+$, the same idea with Hindman's theorem in place of Ramsey's theorem gives the following
generalization.

\begin{proposition}\label{SP:prop:finitegroup}
Let $G$ be a finite abelian group, and let $f\colon\N\to G$ be completely multiplicative, i.e.\
$f(mn)=f(m)f(n)$ for all $m,n\in\N$. Then there are infinite sets $B,C\subset\N$ with
$f(b+c)=f(bc)=1_G$ for all $b\in B$ and $c\in C$.
\end{proposition}

\begin{proof}
Hindman's theorem~\cite[Theorem~3.1]{hindman1974finite}, applied to the finite colouring
$n\mapsto f(n)$, gives $h\in G$ and a sequence $(x_i)_{i\in\N}$ in $\N$ such that $f$ takes the
value $h$ on $\FS((x_i))=\{x_{i_1}+\dots+x_{i_k}:k\in\N,\ i_1<\dots<i_k\}$. We may assume that
$x_1<x_2<\cdots$. Indeed, if $(x_i)$ is unbounded we pass to a strictly increasing subsequence,
whose finite sums form a subset of $\FS((x_i))$. Otherwise some value $x$ occurs infinitely
often, so $\FS((x_i))$ contains every multiple of $x$, and we replace $(x_i)$ by $(2^ix)$. Let $B_0=\{x_1,x_3,\dots\}$ and $C_0=\{x_2,x_4,\dots\}$. Then $f$ equals $h$
on $B_0$, on $C_0$ and on $B_0+C_0$. The image $f(\N)$ is a nonempty subset of the finite group
$G$ that is closed under multiplication, hence a subgroup. So there is $d\in\N$ with
$f(d)=h^{-1}$. Put $B=dB_0$ and $C=dC_0$. For $b_0\in B_0$ and $c_0\in C_0$ we get
$f(db_0+dc_0)=f(d)f(b_0+c_0)=h^{-1}h=1_G$ and
$f(db_0\cdot dc_0)=f(d)^2f(b_0)f(c_0)=h^{-2}h^2=1_G$.
\end{proof}

For instance, for every $k\ge2$ the set $\{n:\Omega(n)\equiv0\pmod k\}$ contains $B+C$ and $BC$
for some infinite $B,C\subset\N$.

\begin{remark}\label{SP:rem:liouville}
\Cref{SP:cor:beatty,SP:thm:liouville} answer Question~\S5.1(a) for both examples named in the
survey. Question~5.2 of~\cite{kra2023problems} itself remains open when ``positive density'' means
that the natural density exists and is positive (see also \cref{SP:rem:q52}).
\end{remark}

%-----------------------------------------------------------------------------------------------
\subsection{Sums and shifted products at logarithmic scale}
\label{SP:sec:q54}

Logarithmic density is defined in~\cite[Definition~3.22]{kra2023problems}: for $A\subset\N$,
\[
\udlog(A)=\limsup_{N\to\infty}\frac1{\log N}\sum_{n=1}^N\frac{\one_A(n)}{n},
\]
with the lower version defined by $\liminf$. When the two agree, their common value is the
logarithmic density $\dlog(A)$. Question~5.4 of~\cite{kra2023problems} reads as follows.

\begin{question}{5.4}\label{SP:q:54}
Suppose $A\subset\N$ has positive upper logarithmic density. Are there infinite sets
$B,C\subset\N$ with $B+C\subset A$ and $BC+1\subset A$?
\end{question}

The survey introduces the shift $BC+1$ because the unshifted question fails for congruence
(local) reasons: the odd numbers contain no $B+C$ together with $BC$, since $b+c$ and $bc$ cannot
both be odd. The survey also notes that Question~5.4 generalizes an earlier question of Moreira
about singletons $B=\{x\}$, $C=\{y\}$. Taken literally, the singleton version is trivial ($x=1$,
$y=a-1$ with $a\in A$); with $x,y\ge2$ it is Moreira's conjecture, which Richter proved for every
set of positive upper logarithmic density~\cite[Theorem~1.5 with $Q(y)=y-1$]{richter2025sums}.
Indeed, that theorem gives $x,y\ge2$ with $x+y-1\in A-1$ and $xy\in A-1$, i.e.\ $x+y\in A$ and
$xy+1\in A$, and $A-1$ has the same upper logarithmic density as $A$. Richter also relates his
results to Questions~5.2 and~5.4 of~\cite{kra2023problems}. We show that the infinite version
fails, because of an obstruction of a different kind: an archimedean one. Our set is aperiodic in
the sense of \eqref{SP:eq:aperiodic}, so it has no congruence obstruction.

Throughout this subsection let $I=[\tfrac{7}{20},\tfrac{13}{20})$, let
$\psi(n)=\log_2 n\bmod 1\in\T$ for $n\in\N$, and let
\[
A=\{n\in\N:\{\log_2n\}\in I\}=\bigcup_{k\ge0}\big[2^{k+7/20},\,2^{k+13/20}\big)\cap\N .
\]
Membership is decided in exact arithmetic: $n\in A$ if and only if
$2^{20k+7}\le n^{20}<2^{20k+13}$, where $k=\lfloor\log_2n\rfloor$.

\begin{proposition}\label{SP:prop:Aprops}
The set $A$ has the following properties.
\begin{enumerate}[label=\textup{(\roman*)}]
\item Its logarithmic density exists and $\dlog(A)=3/10$.
\item It is aperiodic, i.e.\ it satisfies \eqref{SP:eq:aperiodic}.
\item $\udens(A)=2-2^{7/10}\approx0.3755$ and $\ldens(A)=2^{3/10}-1\approx0.2311$. In particular,
the natural density of $A$ does not exist.
\end{enumerate}
\end{proposition}

\begin{proof}
Write $a=7/20$ and $b=13/20$. For real $1\le u<v$, comparing with $\int_u^v dx/x$ gives
$\sum_{u\le n<v}1/n=\log(v/u)+O(1/u)$. So the $k$-th block contributes
$(b-a)\log2+O(2^{-k})$ to $\sum\one_A(n)/n$.

(i) Let $2^K\le N<2^{K+1}$. Summing over the blocks with $k<K$, and bounding the contribution of
the block $k=K$ by $(b-a)\log 2+O(2^{-K})$, we get $\sum_{n\le N}\one_A(n)/n=\tfrac{3}{10}K\log2+O(1)$.
Since $\log N=K\log 2+O(1)$, the claim follows.

(ii) Let $\alpha\in\Q\setminus\Z$. For integers $m\le m'$ we have
$|\sum_{n=m}^{m'}\e(\alpha n)|\le 1/|\sin\pi\alpha|$. The set $A\cap[1,N]$ is a union of at most
$\log_2N+1$ intervals of consecutive integers. Hence
$|\frac1N\sum_{n\le N}\one_A(n)\e(\alpha n)|\le(\log_2N+1)/(N|\sin\pi\alpha|)\to0$.

(iii) Let $N=2^{K+s}$ with $K\in\N$ and $s\in[0,1)$. Counting the blocks with $k<K$ and the part
of block $K$ below $N$ gives
\[
|A\cap[1,N]|=2^K\big(2^b-2^a+g(s)\big)+O(K),\qquad g(s)=\max\big(0,\min(2^s,2^b)-2^a\big).
\]
So $|A\cap[1,N]|/N=\rho(s)+o(1)$, uniformly in $s$, where $\rho(s)=(2^b-2^a+g(s))/2^s$.
On $[0,a]$ the function $\rho(s)=(2^b-2^a)2^{-s}$ decreases. On $[a,b]$ it equals
$1-(2^{a+1}-2^b)2^{-s}$, which increases because $2^{a+1}>2^b$. On $[b,1)$ it equals
$2(2^b-2^a)2^{-s}$, which decreases. Therefore $\min\rho=\rho(a)=2^{b-a}-1=2^{3/10}-1$ and
$\sup\rho=\rho(b)=2-2^{a+1-b}=2-2^{7/10}$. Taking $N=2^{K+a}$ and $N=2^{K+b}$ (rounded down) as
$K\to\infty$ shows that these values are $\ldens(A)$ and $\udens(A)$.
\end{proof}

Let $\eps=\log_2(1+1/58)$. Since $\log(1+x)\le x$ and $\log2>0.693$, we have
$\eps\le1/(58\log2)<0.0249$. The constant $58$ is the least $M\in\N$ with
$2\log_2(1+1/M)<1/20$; this is the margin needed in \cref{SP:thm:q54}(a).

\begin{lemma}\label{SP:lem:shiftphase}
If $x,y\in\N$ satisfy $y\ge58x$ and $x+y\in A$, then
$\{\log_2y\}\in[\tfrac7{20}-\eps,\tfrac{13}{20})$.
\end{lemma}

\begin{proof}
We have $\log_2(x+y)=\log_2y+\eps'$ with $0<\eps'=\log_2(1+x/y)\le\eps<\tfrac7{20}$. Write
$\log_2(x+y)=k+\theta$ with $k\in\Z$ and $\theta=\{\log_2(x+y)\}\in I$. Then
$\log_2y=k+(\theta-\eps')$, and $0\le\theta-\eps'<1$ because $\theta\ge\tfrac7{20}>\eps'$. Hence
$\{\log_2y\}=\theta-\eps'\in[\tfrac7{20}-\eps,\tfrac{13}{20})$.
\end{proof}

\begin{theorem}\label{SP:thm:q54}
\leavevmode
\begin{enumerate}[label=\textup{(\alph*)}]
\item There are no $b_0,c_0,b,c\in\N$ and $\tau\in\{0,1\}$ with $b\ge58c_0$ and $c\ge58b_0$
such that
\[
b_0+c\in A,\qquad c_0+b\in A,\qquad bc+\tau\in A .
\]
\item Let $t_1,t_2\in\Z$. There are no infinite sets $B,C\subset\N$ with $b+c+t_1\in A$ and
$bc+t_2\in A$ for all $b\in B$ and $c\in C$.
\end{enumerate}
In particular, \cref{SP:q:54} has a negative answer: $\dlog(A)=3/10$, but there are no infinite
sets $B,C\subset\N$ with $B+C\subset A$ and $BC+1\subset A$.
\end{theorem}

\begin{proof}
(a) By \cref{SP:lem:shiftphase}, applied to $(x,y)=(b_0,c)$ and to $(x,y)=(c_0,b)$,
\[
\{\log_2c\},\ \{\log_2b\}\in\big[\tfrac7{20}-\eps,\ \tfrac{13}{20}\big).
\]
Write $\log_2(bc+\tau)=\log_2b+\log_2c+\eta$, where $\eta=\log_2(1+\tau/(bc))$. Since
$bc\ge58^2$, we have $0\le\eta\le1/(58^2\log2)<0.001$. Let $u=\{\log_2b\}+\{\log_2c\}+\eta$.
Then $\log_2(bc+\tau)-u\in\Z$, so $\{\log_2(bc+\tau)\}=\{u\}$, and
$u\in[\tfrac{7}{10}-2\eps,\ \tfrac{13}{10}+\eta)$.
If $u<1$, then $\{u\}=u\ge\tfrac7{10}-2\eps>\tfrac{13}{20}$, because $2\eps<0.05$.
If $u\ge1$, then $\{u\}=u-1<\tfrac3{10}+\eta<\tfrac7{20}$.
In both cases $\{\log_2(bc+\tau)\}\notin I$, so $bc+\tau\notin A$.

(b) Suppose such $B$ and $C$ exist. Fix $b_0\in B$. For $c\in C$, $c\to\infty$, the distance in
$\T$ between $\psi(c+b_0+t_1)$ and $\psi(c)$ is at most $|\log_2(1+(b_0+t_1)/c)|$, which tends to $0$.
Since $\psi(c+b_0+t_1)\in I$, every limit point of $\psi(c)$ as $c\to\infty$ in $C$ lies in the
closure $\bar I=[\tfrac7{20},\tfrac{13}{20}]$. By symmetry the same holds for $\psi(b)$ as $b\to\infty$ in $B$.
Choose $b_n\in B$ and $c_n\in C$ with $b_n,c_n\to\infty$. Since $\T$ is compact, after passing
to a subsequence we may assume that $\psi(b_n)\to\beta$ and $\psi(c_n)\to\gamma$, where
$\beta,\gamma\in\bar I$. Now
$\psi(b_nc_n+t_2)=\psi(b_n)+\psi(c_n)+\log_2(1+t_2/(b_nc_n))\to\beta+\gamma$ in $\T$. But
$\beta+\gamma\in[\tfrac7{10},\tfrac{13}{10}]$, and modulo $1$ this set is
$[\tfrac7{10},1)\cup[0,\tfrac3{10}]$, which is at distance $\tfrac1{20}$ from $\bar I$. On the other hand
$\psi(b_nc_n+t_2)\in I$ for all $n$, so $\beta+\gamma\in\bar I$. This is a contradiction.

For the final statement, recall $\dlog(A)=3/10$ from \cref{SP:prop:Aprops}(i). If $B,C\subset\N$
were infinite with $B+C\subset A$ and $BC+1\subset A$, we could fix $b_0\in B$ and $c_0\in C$ and
then pick $b\in B$ with $b\ge58c_0$ and $c\in C$ with $c\ge58b_0$, contradicting (a) with
$\tau=1$. (It also follows from (b) with $t_1=0$, $t_2=1$.)
\end{proof}

\begin{remark}\label{SP:rem:q54}
(i) The obstruction is asymptotic, and it needs both $B$ and $C$ to be infinite. If $B=\{1\}$, then
$B+C=BC+1=C+1$, so $C=(A-1)\cap\N$ works. By Richter's theorem quoted above, every set of positive
upper logarithmic density, in particular $A$, contains $\{b+c,bc+1\}$ with $b,c\ge2$. Larger finite
configurations also occur in $A$: for $B=\{1,2,4,8,16\}$ and $C=\{1306,1307,\dots,1310\}$, all
$25$ sums $b+c$ and all $25$ numbers $bc+1$ lie in $A$. This is a direct check of $50$ memberships
with the exact test above. Part~(a) of \cref{SP:thm:q54} is a quantitative form of the
obstruction: it rules out a four-element configuration in which $b,c$ are large compared with
$c_0,b_0$.

(ii) By \cref{SP:thm:q54}(b), shifting the sums or the products by any fixed integers does not
remove the obstruction. Unlike the congruence obstruction to the unshifted question exhibited by the odd numbers,
it is not local: $A$ is aperiodic (\cref{SP:prop:Aprops}(ii)). It is also thick (the blocks have
unbounded length), so $\ubd(A)=1$. Similar constructions at logarithmic scale appear throughout
\S\S2--3 of the survey~\cite{kra2023problems}, where they serve as counterexamples for sumsets
$B+B+t$; Richter's Example~1.6 in~\cite{richter2025sums} is another block-type construction of
this kind, a set of positive upper density without $\{x+y-1,xy\}$, $x,y\ge2$. Here the same idea is applied
to infinite sums and shifted products. To our knowledge this counterexample is new.

(iii) The set $A$ is not syndetic. The survey also asks Question~5.4 for syndetic $A$, and
Question~5.5 of~\cite{kra2023problems} (whether every syndetic set contains $BC$ with $B,C$ infinite);
both remain open.
\end{remark}

Recall Question~5.2 of~\cite{kra2023problems}.

\begin{question}{5.2}\label{SP:q:52}
Fix $A\subset\N$ that has positive density and is aperiodic. Are there infinite sets
$B,C\subset\N$ with $B+C\subset A$ and $BC\subset A$?
\end{question}

\begin{remark}\label{SP:rem:q52}
Taking $\tau=0$ in \cref{SP:thm:q54}(a), or $t_1=t_2=0$ in (b), shows that $A$ contains no $B+C$
together with $BC$ for infinite $B,C\subset\N$. By \cref{SP:prop:Aprops}, $A$ is aperiodic and has
positive upper density (and positive logarithmic and upper Banach density), but its natural density
does not exist. So $A$ gives a negative answer to \cref{SP:q:52} if ``positive density'' there is
read as positive upper density (or positive upper Banach or logarithmic density). This reading is
consistent with the survey's text, which says that (19) is equivalent to equidistribution in
residue classes ``if the density of $A$ exists''. It does not settle \cref{SP:q:52} when the
natural density is required to exist; that case remains open.
\end{remark}

%-----------------------------------------------------------------------------------------------
\subsection{Sums and products in \texorpdfstring{$\Q$}{Q}}
\label{SP:sec:q56}

Following~\cite[\S5.1]{kra2023problems}, where the notion is attributed
to~\cite[Definition~1.3]{bergelson2017ergodic}, we use the formulation
of~\cite[Definition~1.7]{kousek2024revisiting}: a \emph{double \Fol{} sequence} in $\Q$ is a
sequence $\Phi=(F_N)_{N\in\N}$ of nonempty finite subsets of $\Q$ such that
\[
\lim_{N\to\infty}\frac{|F_N\cap(x+F_N)|}{|F_N|}=1\quad(x\in\Q)
\qquad\text{and}\qquad
\lim_{N\to\infty}\frac{|F_N\cap xF_N|}{|F_N|}=1\quad(x\in\Q^*).
\]
The survey describes $\Phi$ as simultaneously a \Fol{} sequence in $(\Q,+)$ and in
$(\Q^*,\times)$; strictly speaking the latter requires $F_N\subset\Q^*$. This makes no
difference below: $0$ is a single point and $|F_N|\to\infty$ (\cref{SP:lem:folner}(i)), so all
densities are unchanged if $0$ is removed from every $F_N$.
For $A\subset\Q$ we write $\udens_\Phi(A)=\limsup_N|A\cap F_N|/|F_N|$. When the limit exists we
denote it by $\dens_\Phi(A)$, as in~\cite[\S1]{kra2023problems}.

\begin{question}{5.6}\label{SP:q:56}
Let $A\subset\Q$ have $\dens_\Phi(A)>0$ for some double \Fol{} sequence $\Phi$. Are there infinite
sets $B,C\subset\Q$ such that $B+C\subset A$ and $BC\subset A$?
\end{question}

We show that the answer is negative, for every double \Fol{} sequence.

\begin{lemma}\label{SP:lem:folner}
Let $\Phi=(F_N)$ be a double \Fol{} sequence in $\Q$.
\begin{enumerate}[label=\textup{(\roman*)}]
\item $|F_N|\to\infty$.
\item For every bounded set $X\subset\Q$ we have $|F_N\cap X|/|F_N|\to0$.
\item $|F_N\cap\Q_{<0}|/|F_N|\to1/2$.
\end{enumerate}
\end{lemma}

\begin{proof}
(i) Suppose $|F_N|\le M$ for infinitely many $N$. For those $N$ that are large enough,
$|F_N\cap(1+F_N)|>(1-\frac1{M+1})|F_N|\ge|F_N|-1$, hence $F_N\subset1+F_N$. Both sets have the same
cardinality, so $F_N=1+F_N$. This is impossible for a finite nonempty $F_N\subset\Q$, since
$\max(1+F_N)=1+\max F_N>\max F_N$.

(ii) Choose $R\in\N$ with $X\subset[-R,R]$, fix $k\in\N$, and let $t_j=j(2R+1)$ for
$0\le j<k$. The sets $X+t_j$ are pairwise disjoint. For each $t\in\Q$,
\[
|F_N\cap(X+t)|=|(F_N-t)\cap X|\ge|F_N\cap X|-|F_N\setminus(F_N-t)|=|F_N\cap X|-o(|F_N|),
\]
since $|F_N\setminus(F_N-t)|=|F_N|-|F_N\cap(-t+F_N)|$. Summing over $j$ gives
$k|F_N\cap X|\le|F_N|+o(|F_N|)$, so $\limsup_N|F_N\cap X|/|F_N|\le1/k$. Now let $k\to\infty$.

(iii) Let $p_N=|F_N\cap\Q_{>0}|$ and $n_N=|F_N\cap\Q_{<0}|$. The map $y\mapsto-y$ gives
$n_N=|(-F_N)\cap\Q_{>0}|$, and this differs from $p_N$ by at most
$|(-F_N)\setminus F_N|=|F_N|-|F_N\cap(-1)F_N|=o(|F_N|)$. Moreover $p_N+n_N\ge|F_N|-1$. By (i),
$n_N/|F_N|\to1/2$.
\end{proof}

\begin{lemma}\label{SP:lem:noconf}
Let $A\subset\Q_{<0}$ be such that $A\cap[-R,R]$ is finite for every $R>0$. Then there are no
infinite sets $B,C\subset\Q$ with $B+C\subset A$ and $BC\subset A$.
\end{lemma}

\begin{proof}
Suppose such $B,C$ exist. Every product $bc$ is negative, so $0\notin B\cup C$, and every
$b\in B$ and $c\in C$ have opposite signs. If $B$ contained a positive and a negative element,
no $c\in C$ could have the sign opposite to both. So $B$ has constant sign, and $C$ has the
opposite constant sign. By symmetry of the hypotheses in $B$ and $C$ we may assume that
$B\subset\Q_{<0}$ and $C\subset\Q_{>0}$. Fix $b^*\in B$. For $c\in C$ we have $b^*+c<0$, so
$0<c<|b^*|$ and $b^*c\in(-(b^*)^2,0)$. Since $c\mapsto b^*c$ is injective and $C$ is infinite,
$b^*C$ is an infinite subset of $A\cap[-(b^*)^2,(b^*)^2]$, which is a contradiction.
\end{proof}

\begin{theorem}\label{SP:thm:q56}
Let $\Phi=(F_N)$ be any double \Fol{} sequence in $\Q$. There is a set $A\subset\Q_{<0}$ such
that $\dens_\Phi(A)=\lim_{N}|A\cap F_N|/|F_N|=1/2$, and such that there are no infinite sets
$B,C\subset\Q$ with $B+C\subset A$ and $BC\subset A$. In particular, \cref{SP:q:56} has a
negative answer.
\end{theorem}

\begin{proof}
For $R\in\N$ write $[-R,R]$ for $\Q\cap[-R,R]$, and let $\eps_R(N)=|F_N\cap[-R,R]|/|F_N|$. By
\cref{SP:lem:folner}(ii), $\eps_R(N)\to0$ as $N\to\infty$ for each $R$. So we can choose
$N_1<N_2<\cdots$ such that $\eps_R(N)<1/R$ for all $N\ge N_R$. Put
\[
K_R=\bigcup_{N<N_{R+1}}F_N\ \ (\text{a finite set}),\qquad
Z=\bigcup_{R\in\N}\big([-R,R]\setminus K_R\big),\qquad A=\Q_{<0}\setminus Z .
\]
\emph{Bounded parts of $A$ are finite.} Let $R\in\N$. If $x\in A$ and $|x|\le R$, then
$x\notin[-R,R]\setminus K_R$, so $x\in K_R$. Hence $A\cap[-R,R]\subset K_R$ is finite. By
\cref{SP:lem:noconf}, $A$ contains no $B+C$ together with $BC$ for infinite $B,C$.

\emph{$Z$ has density zero.} Let $N\ge N_1$, and let $R$ be such that $N_R\le N<N_{R+1}$.
If $R'\ge R$, then $N<N_{R'+1}$, so $F_N\subset K_{R'}$ and $F_N$ does not meet
$[-R',R']\setminus K_{R'}$. If $R'<R$, then $[-R',R']\subset[-R,R]$. Hence
$F_N\cap Z\subset F_N\cap[-R,R]$, and so $|F_N\cap Z|\le\eps_R(N)|F_N|<|F_N|/R$. As
$N\to\infty$, also $R\to\infty$, so $|F_N\cap Z|/|F_N|\to0$.

\emph{Density of $A$.} By \cref{SP:lem:folner}(iii),
$|A\cap F_N|/|F_N|=|F_N\cap\Q_{<0}|/|F_N|-O(|F_N\cap Z|/|F_N|)\to1/2$.
\end{proof}

\begin{remark}\label{SP:rem:q56easy}
We record some easy facts that put \cref{SP:thm:q56} in context. Let $\Phi$ be a double \Fol{}
sequence and let $A\subset\Q$ satisfy $\udens_\Phi(A)>0$.
\begin{enumerate}[label=\textup{(\roman*)}]
\item Each of the two conditions can be met separately. Moreira, Richter and Robertson proved the
following~\cite[Theorem~1.3]{moreira2019proof}: if $G$ is a countable group, $\Psi$ is a two-sided
\Fol{} sequence on $G$, and $E\subset G$ has positive upper density along $\Psi$, then $E\supset BC$
for some infinite $B,C\subset G$. Applied to $G=(\Q,+)$ and $\Psi=\Phi$, this gives infinite $B,C$
with $B+C\subset A$. Next apply it to $G=(\Q^*,\cdot)$ and $\Psi=(F_N\setminus\{0\})_{N\ge N_0}$, where $N_0$ is such that these sets are nonempty; this is a
\Fol{} sequence in $\Q^*$ by \cref{SP:lem:folner}(i), and the upper density of $A\setminus\{0\}$
along it equals $\udens_\Phi(A)$. This gives infinite $B,C\subset\Q^*$ with $BC\subset A$. (Both
groups are abelian, so one-sided and two-sided \Fol{} sequences coincide.)
\item The set $\Q_{<0}$ itself contains the configuration: take $B=\{-n:n\ge2\}$ and
$C=\{1/m:m\in\N\}$. So the deletion of the null set $Z$ in \cref{SP:thm:q56} is essential.
\item The question requires two sets $B,C$. The version with $B=C$ fails even for restricted
products: if $B\subset\Q$ is infinite, then $B\setminus\{0\}$ has two distinct elements of the
same sign. So $\{b_1b_2:b_1,b_2\in B,\ b_1\neq b_2\}$ always contains a nonnegative number, and it
is not contained in $\Q_{<0}$. On the other hand, $\dens_\Phi(\Q_{<0})=1/2$ by
\cref{SP:lem:folner}(iii).
\end{enumerate}
\end{remark}

\begin{remark}\label{SP:rem:q56literal}
\Cref{SP:thm:q56} answers \cref{SP:q:56} as stated. The question assumes only that $A$ has positive
density along some double \Fol{} sequence, and our set has density $1/2$ along any prescribed
one. The obstruction is the sign of $x$, i.e.\ the real place of $\Q$. It plays the same role as
congruence obstructions in $\N$, which is why the survey adds the aperiodicity hypothesis to
Question~5.2. The point of \cref{SP:thm:q56} is that, once a set of density zero has been removed,
this obstruction cannot be detected by density along a double \Fol{} sequence. The $\Phi$-dependent deletion of finite pieces is similar to the device used
in~\cite[Example~5.9]{kra2023problems} (compare also the proof
of~\cite[Proposition~3.2]{kra2023problems}). The question that is
presumably intended remains open: does every $A\subset\Q$ that is large and has no such local
obstruction contain $B+C$ and $BC$? Possible formalizations are to require $\ldens_\Phi(A)>0$
for \emph{every} double \Fol{} sequence $\Phi$, which our construction (it depends on $\Phi$) does
not address, or to impose a hypothesis on $A$ analogous to aperiodicity. We do not know the answer
in either case.
\end{remark}
